\uuid{q3RE}
\exo7id{1654}
\titre{exo7 1654}
\auteur{legall}
\organisation{exo7}
\datecreate{1998-09-01}
\isIndication{false}
\isCorrection{false}
\chapitre{Réduction d'endomorphisme, polynôme annulateur}
\sousChapitre{Diagonalisation}
\cpge{mpsi-mp2i.matrices-applications-lineaires.a, pcsi.matrices-applications-lineaires.a, ptsi.matrices-applications-lineaires.a, mpsi-mp2i.changements-bases.a, pcsi.matrices-applications-lineaires.c, ptsi.matrices-applications-lineaires.c, mp-mpi.reduction-endomorphismes.b, pc.reduction.a, psi.reduction-endomorphismes.a, pt.reduction-endomorphismes.a, mp-mpi.reduction-endomorphismes.d, pc.reduction.c, psi.reduction-endomorphismes.c, pt.reduction-endomorphismes.c, mp-mpi.reduction-endomorphismes.i, pc.reduction.d, psi.reduction-endomorphismes.d, bcpst2.valeurs-propres-vecteurs-propres.a, bcpst2.valeurs-propres-vecteurs-propres.b}
\module{Algèbre}
\niveau{L2}
\difficulte{}

\contenu{
\texte{
Soit $  f \in \mathcal{L} (\R ^3)  $ d\'etermin\'ee par sa matrice $  M=\begin{pmatrix}
1 & 1 & 1 \cr
                                      1 & 1 & 1 \cr
                                      -1 & 1 & 1 \cr \end{pmatrix} $ dans une base $  \{ e_1,
e_2,e_3\}   $ de $  \R ^3  .$
}
\begin{enumerate}
    \item \question{Montrer que $  M  $ est diagonalisable.}
    \item \question{Montrer que la restriction de $  f  $ a tout sous-espace stable est diagonalisable.}
    \item \question{En d\'eduire tous les sous-espaces de $  \R ^3  $ stables par $  f  .$}
\end{enumerate}
}
